% Part: normal-modal-logic
% Chapter: syntax-and-semantics
% Section: relational-models

\documentclass[../../../include/open-logic-section]{subfiles}

\begin{document}

\olfileid{nml}{syn}{rel}

\olsection{సంబంధాత్మక నమూనాలు}

నార్మల్ మోడల్ తర్కాల అర్థవిచారంలో
మౌలిక భావన \emph{సంబంధాత్మక నమూనా}.
దానిలో లోకాల సమితి, వాటిని కలిపే
ద్విస్థానిక “ప్రాప్యత సంబంధం,”
ఏ లోకంలో ఏ
\tetoken{ప్రతిజ్ఞావాక్య చరాలు}{!!{propositional variable}s}
“సత్యం”గా పరిగణించాలో నిర్ణయించే
కేటాయింపు ఉంటాయి.

\begin{defn}
  మౌలిక మోడల్ భాషకు ఒక \emph{నమూనా}
  $\mModel{M} = \tuple{W, R, V}$ అనే
  త్రయం. ఇందులో
  \begin{enumerate}
  \item $W$ అనేది “లోకాల” ఖాళీ కాని సమితి,
  \item $R$ అనేది $W$పై ద్విస్థానిక ప్రాప్యత సంబంధం,
  \item $V$ అనేది ప్రతి
    \tetoken{ప్రతిజ్ఞావాక్య చరానికి}{!!{propositional
    variable}}~$p$ సాధ్య లోకాల
    సమితి $V(p)$ని కేటాయించే ప్రమేయం.
  \end{enumerate}
  $Rww'$ నిలిస్తే, $w'$ అనేది
  $w$ నుంచి \emph{ప్రాప్యమైనది}
  అంటాం. $w \in V(p)$ అయితే,
  $p$ అనేది $w$లో \emph{సత్యం}
  అంటాం.
\end{defn}

సంబంధాత్మక అర్థవిచారపు గొప్ప సౌలభ్యం:
\olref{fig:simple}లో ఉన్నటువంటి
సరళ చిత్రాలతో నమూనాలను చూపవచ్చు.
లోకాలను బిందువులు సూచిస్తాయి;
$w$ నుంచి $w'$కు బాణం ఉన్నప్పుడు,
అప్పుడు మాత్రమే $w'$ అనేది
$w$ నుంచి ప్రాప్యమైనది. అంతేకాక,
$w \in V(p)$ అయితే బిందువుకు
(లోకానికి)~$\mTrue{p}$ గుర్తు,
లేకపోతే~\mFalse{p} గుర్తు పెడతాం.
\olref{fig:simple} చిత్రం
$W = \{w_1, w_2, w_3\}$,
$R = \{\tuple{w_1, w_2},\tuple{w_1,
  w_3}\}$, $V(p) = \{w_1, w_2\}$,
మరియు $V(q) = \{w_2\}$ ఉన్న
నమూనాను చూపుతుంది.

\begin{figure}
  \begin{center}
    \begin{tikzpicture}[modal]
      \node[world] (w1) [label={[align=right]right:\mTrue{p}\\ \mFalse{q}}]{$w_1$};
      \node[world] (w2) [label={[align=right]right:\mTrue{p}\\ \mTrue{q}},
        above right=of w1]{$w_2$};
      \node[world] (w3) [label={[align=right]right:\mFalse{p}\\ \mFalse{q}},
        below right=of w1] {$w_3$};
      \draw[->] (w1) to (w2);
      \draw[->] (w1) to (w3);
    \end{tikzpicture}
  \end{center}
  \caption{ఒక సరళ నమూనా.}
  \ollabel{fig:simple}
\end{figure}

\end{document}
